The star $\alpha$-Centauri (Rigel Kentaurus) is a triple star which consists
of two main-sequence stars $\alpha$-Centauri A and $\alpha$-Centauri B
representing a visual binary system, and the third star, called Proxima
Centauri, which is smaller and fainter than the other two stars. The angular
distance between $\alpha$-Centauri A and $\alpha$-Centauri B is $17.59''$.
The binary system has an orbital period of 79.24 years. The visual magnitudes
of $\alpha$-Centauri A and $\alpha$-Centauri B are $-0.01$ and 1.34
respectively. Their color indices are 0.65 and 0.85 respectively. Use the
data below to answer the following questions.

\begin{center}
Data for main-sequence stars

\medskip
\begin{tabular}{ccc}
$(B-V)_0$ & $T_{\mathrm{eff}}$ & BC \\
$-0.25$ & 24500 & 2.30 \\
$-0.23$ & 21000 & 2.15 \\
$-0.20$ & 17700 & 1.80 \\
$-0.15$ & 14000 & 1.20 \\
$-0.10$ & 11800 & 0.61 \\
$-0.05$ & 10500 & 0.33 \\
0.00 & 9480 & 0.15 \\
0.10 & 8530 & 0.04 \\
0.20 & 7910 & 0 \\
0.30 & 7450 & 0 \\
0.40 & 6800 & 0 \\
0.50 & 6310 & 0.03 \\
0.60 & 5910 & 0.07 \\
0.70 & 5540 & 0.12 \\
0.80 & 5330 & 0.19 \\
0.90 & 5090 & 0.28 \\
1.00 & 4840 & 0.40 \\
1.20 & 4350 & 0.75 \\
\end{tabular}

\medskip
BC = Bolometric Correction, $(B-V)_0$ = Intrinsic Color
\end{center}

Questions:
\begin{enumerate}
\item Plot the curve BC versus $(B-V)_0$.
\item Determine the apparent bolometric magnitudes of $\alpha$-Centauri A
  and $\alpha$-Centauri B using the corresponding curve.
\item Calculate the mass of each star.
\end{enumerate}

Notes:
\begin{enumerate}
\item Bolometric correction (BC) is a correction that must be made to the
  apparent magnitude of an object in order to convert an object's visible
  magnitude to its bolometric magnitude:
  \[ \mathrm{BC} = m_v - m_{\mathrm{bol}} \quad \mbox{or} \quad
     \mathrm{BC} = M_v - M_{\mathrm{bol}} \]
\item Luminosity mass relation:
  \[ M_{\mathrm{bol}} = -10.2\log\left(\frac{M}{M_\odot}\right) + 4.9 \]
\end{enumerate}
